\section{Introduction}

Schizophrenia is a devastating psychiatric disorder that appears to result in
part from the functional disruption of prefrontal cortical networks
\citep{insel2010rethinking, jauhar2022schizophrenia}. Despite decades of research,
a central gap remains: we lack a mechanistic understanding of how molecular and
synaptic disruptions in prefrontal circuits cause the network-level failures
observed in the disease. Closing this gap is critical because it would explain
how the genetic and environmental risks that alter synaptic function cause
prefrontal networks to fail.

Parallel nonhuman primate and mouse genetic models of prefrontal network failure
have begun to characterize this link at the physiological level
\citep{zick2018blocking, kummerfeld2020cognitive, zick2021disparate}. Diverse
insults operating through distinct molecular mechanisms---blocking NMDA
receptors (NMDAR) in primates with phencyclidine, or deleting a schizophrenia
risk gene in mice related to microRNA processing---converge on the same
endpoint in prefrontal cortex: reduced 0-lag synchronous spiking between
prefrontal neurons and reduced information transfer at synaptic interaction
lags \citep{zick2018blocking, zick2021disparate}. These convergent
pathophysiological signatures suggest that multiple risk factors act on a
common vulnerability of prefrontal local circuits. However, these observations
do not explain \textit{how} or \textit{why} molecular disruptions of synaptic
function produce the observed distortions of network dynamics.

To address this question, we develop a spiking neural network model that is
informed by, and provides mechanistic explanation for, neural recordings
obtained in monkey prefrontal cortex before and during NMDAR blockade
\citep{zick2018blocking}. We use mean field analysis and linear stability
analysis to investigate how the balance between AMPA and NMDA components of
recurrent excitation and GABA inhibition shapes regimes of network dynamics
and spiking synchrony. Theoretical studies have established that synchrony can
arise in randomly connected recurrent networks operating in the asynchronous
irregular regime \citep{brunel2000dynamics, vanvreeswijk1996chaos,
renart2010asynchronous} and the synchronous irregular regime
\citep{brunel2000dynamics, brunel1999fast, brunel2003fast}. In both regimes,
individual neurons fire highly irregularly at low rates, as is typical of
cortex. The critical distinction is that in the asynchronous regime, population
spike rate is steady in time, whereas in the synchronous regime it becomes
oscillatory.

We show that a simulated prefrontal network operating near the boundary between
these two dynamical regimes can explain several key experimental observations.
First, such a network achieves biologically realistic stochastic spike trains
and firing rates for excitatory and inhibitory neurons. Second, a modest
increase in extrinsic inputs---such as might occur during behavior---shifts
the critical network from a steady to an oscillatory regime, producing emergent
0-lag spike synchrony as neurons stochastically entrain to oscillatory
population activity. Third, and most importantly, we show that reducing
recurrent NMDAR synaptic currents prevents this transition, thereby preventing
the emergence of 0-lag spike synchrony. This establishes direct parallels with
the biological data, including the emergence of synchronous spiking via
recurrent synaptic interactions during behavior, its association with gamma
oscillation, and the joint dependence of both on NMDAR synaptic mechanisms
\citep{zick2018blocking}.

\paragraph{Contributions.}
\begin{itemize}
  \item We construct a spiking network model of prefrontal local circuits that
    operates near the boundary between asynchronous irregular and synchronous
    irregular regimes.
  \item We use mean field and linear stability analysis to derive network
    parameters that place the system at this critical boundary with
    biologically realistic firing rates.
  \item We demonstrate that small increases in extrinsic input shift the
    critical network into an oscillatory regime with emergent 0-lag spike
    synchrony.
  \item We show that reducing recurrent NMDAR currents prevents this
    transition, providing a circuit-level mechanism for the desynchronization
    observed in primate PFC following NMDAR blockade \citep{zick2018blocking}.
  \item We link NMDAR synaptic mechanisms, oscillatory dynamics, and spike
    timing to a potential active disconnection mechanism in schizophrenia
    \citep{dan2004spike, zick2021disparate}.
\end{itemize}

\section{Background and Related Work}

\paragraph{Balanced networks and irregular cortical dynamics.}
The theoretical framework of balanced networks, established by
\citet{vanvreeswijk1996chaos}, showed that cortical-like irregular spiking can
emerge in large networks of excitatory and inhibitory neurons connected sparsely
by strong synapses. \citet{brunel2000dynamics} extended this to a rich
repertoire of dynamical states, including asynchronous irregular and
synchronous irregular regimes, and characterized transitions between them.
\citet{renart2010asynchronous} further demonstrated that such asynchronous
states naturally arise in recurrent cortical circuits. The mean field
framework for analyzing these networks was developed by
\citet{amit1997model, brunel2000dynamics} and extended to networks with
conductance-based synapses by \citet{brunel2001effects}.

\paragraph{Network oscillations and synaptic dynamics.}
\citet{brunel2003fast} established that the balance between recurrent AMPA
excitation and GABA inhibition determines the frequency of fast network
oscillations, while \citet{brunel1999fast} analyzed the emergence of fast
global oscillations in networks of integrate-and-fire neurons with low firing
rates. These studies provide the theoretical basis for understanding how
changes in synaptic balance can shift a network between dynamical regimes.

\paragraph{NMDAR modulation in cortical network models.}
The impact of NMDAR synaptic transmission on neural dynamics has been
simulated in various contexts: the stability of working memory
\citep{wang1999synaptic, compte2000synaptic}, LFP oscillations
\citep{brunel2003fast}, and neuromodulation of working memory circuits
\citep{brunel2001effects}. However, these studies did not address how NMDAR
hypofunction desynchronizes spiking activity in prefrontal cortex or how
spiking synchrony is modulated during transitions between oscillatory and
steady states. By presenting evidence that prefrontal networks operate near
the boundary between these regimes, our model provides a unique circuit-level
insight explaining why NMDAR synaptic deficits reduce both oscillations and
synchronous spiking in prefrontal networks, as observed experimentally
\citep{zick2018blocking}.

\paragraph{Spike timing, plasticity, and schizophrenia.}
Synchronous spiking governs spike-timing-dependent plasticity
\citep{dan2004spike}, establishing a direct link between network dynamics
and synaptic connectivity. Recent evidence further suggests that synchronous
pre- and postsynaptic spiking establishes eligibility traces that render
specific dendritic spines susceptible to potentiation by subsequent dopamine
signals \citep{he2015eligibility, yagishita2014critical}. Any insult that
disrupts the balance between NMDAR inputs, oscillatory activity, and spike
timing in prefrontal circuits could therefore cause these circuits to
actively disconnect via spike-timing-dependent synaptic plasticity, a
hypothesis advanced by \citet{zick2018blocking, zick2021disparate}.
